\documentclass{article}
\usepackage{amsmath}
\begin{document}
\section*{Frame freshness and the option to carry a preexisting local field}
A command can preserve its earlier local vertical field and still
violate the old world-vertical requirement at consumption. This
study links the shared-slot continuation problem to a consumer-frame
uncertainty model. It is a bound on the selected setter input, not a
claim that a body physically stopped or reached a pose.

\paragraph{Declared command interface.}
At the last qualified frame observation the reference axes align.
By the consume cut, pitch may be any $\theta\in[-\rho,\rho]$
with $0\le\rho<\pi/2$. The commanded local pair $(x,z)$ is
transformed to $(x\cos\theta+z\sin\theta,\,
z\cos\theta-x\sin\theta)$. An already live duty requires the
world-vertical effect to lie in $[\ell,u]$ with $0<\ell\le u$;
the horizontal effect must have absolute value at most $\delta$.
There is only one frozen command at the decision cut. In the
carry interface the verifier may change the local $x$ coordinate
but keeps the preexisting $z_0>0$. In the composed interface an
explicitly authorized writer may choose both.
These are different declared privileges. In particular the local
$x$ writer is not by that fact an authorized, strictly
world-horizontal contribution writer. The model compares aggregate
effects; it does not infer separate world provenance from one setter.

\paragraph{Sharp bounds.}
For any feasible $(x,z)$, $( -x,z)$ is feasible by the symmetry of
the frame set and the horizontal tolerance. Intersection across all
frames is convex, so $x=0$ can be chosen without loss of feasibility.
For $x=0$, the vertical range runs from $z\cos\rho$ to $z$ and
the maximum absolute horizontal effect is $z\sin\rho$. It follows
that the carry interface is feasible exactly when
\[
z_0\cos\rho\ge \ell,\qquad z_0\le u,\qquad z_0\sin\rho\le\delta .
\]
The composed interface is feasible exactly when
\[
\frac{\ell}{\cos\rho}\le u,\qquad \ell\tan\rho\le\delta ,
\]
attained by $z=\ell/\cos\rho$. In particular, reconstructing the
actual effects at a sampled older pose is not a test of a later
consume cut. If the source supplies a valid symmetric error and
pitch-drift bound $\rho\le\rho_0+\omega a$ at evidence age $a$,
these give a maximum usable age under that guarantee. The
bounds are necessary for the declared full interval, not a
claim all angles of a real source are attainable. A different
observation law or restrictions on writable $x$ and $z$ need
their own feasible-set check.

\paragraph{Discriminating source controls.}
For $\rho=\arctan(1/2)$, $\ell=0.6$, $u=0.7$,
$\delta=0.31$, and staged $z_0=5/8$, a carried vertical field
cannot certify the interval. The authorized combined command
can. At a known positive extreme the same local-$x$-only
interface can choose $x=-z_0\tan\rho$ and satisfy both world
bounds, but that value does not work at the opposite extreme.
A separate candidate passes the two extreme pitches and violates
the upper vertical bound at the center. Checking only the two
endpoints is not a valid certificate for a prescribed command.

The pinned native velocity writer and physics setter are executed
with an authored pitch-only rotation stand-in. Each case starts
from the native lift's earlier staged component. The two-action
interfaces denote choices entering the same vector write, not
permissions provided by that source. Twelve native setter cases
and 384 independent finite-scenario LP checks agree with the
bounds, with the continuum property supplied by the proof.
There is no sampled frame freshness guarantee, hardware motion
claim, matched advantage over a direct guard with the same
privileges, or independent mission result.
\end{document}
